\documentclass[11pt]{article}
\usepackage[margin=1in]{geometry}
\usepackage{amsmath,amssymb,amsthm}
\usepackage{booktabs,tabularx,array}
\usepackage{hyperref}
\usepackage{xcolor}
\usepackage{enumitem}
\usepackage{placeins}
\hypersetup{colorlinks=true,linkcolor=blue!50!black,urlcolor=blue!50!black,citecolor=blue!50!black}
\newcommand{\code}[1]{\texttt{#1}}
\newcolumntype{L}{>{\raggedright\arraybackslash}X}

\title{Frozen, Certified, and Often Wrong:\\
Lessons from Formalizing 114 Open Problems in Information Theory in Lean~4}
\author{The \emph{Open Problems in Information Theory} contributors\thanks{%
Repository: \url{https://github.com/gokhanmergen/information-theory-problems}; site:
\url{https://gokhanmergen.github.io/information-theory-problems/}. Every number in this note is
reproducible from the repository at the commits cited. \textbf{DRAFT}: prepared with the help of AI
agents. The author list, venue, and final framing are for the maintainers to decide.}}
\date{Draft of 27 November 2026}

\emergencystretch=3em
\begin{document}
\maketitle

\begin{abstract}
We describe a public catalog of 114 open problems in information theory in which every
problem is paired with a \emph{frozen} Lean~4/Mathlib proposition \code{Statement : Prop}.
Anyone, human or machine, may settle a problem by giving a kernel-checked proof of
\code{Statement} or of \code{\textasciitilde Statement} that uses only Lean's standard axioms.
Such a benchmark is only as meaningful as its formal statements. A proof of a mis-formalized
proposition settles nothing, and a mis-formalized proposition is exactly what an automated prover
will find first.
We therefore attach to every statement a machine-checked \emph{certificate}: Lean proofs that the
exact objects in \code{Statement} satisfy the known bounds quoted in the problem's docstring, and
regression tests against textbook special cases. These are backed by a 70k-line library of
formalized Shannon theory.
Building the certificates and auditing provenance changed a large share of the catalog.
Of the 109 statements in the first frozen release, 24 (22\%) had to be semantically repaired, some
more than once. Two more turned out to be theorems already in the literature, and one, a
formulation of the catalog's own, was formally disproved. We classify these failures into six
types, describe the checks that caught each, and give recommendations for anyone freezing formal
statements of open problems.
We also report the catalog's first formal results: a kernel-checked disproof, and a
machine-checked bridge from an independent published Lean proof of the Courtade--Kumar conjecture
to the catalog's frozen statement. The bridge proves that two independently written formalizations
of $I(f(X);Y)$ are equal.
\end{abstract}

\section{Introduction}

Lists of open problems shape research. Erd\H{o}s's problems, now catalogued online, show that a
precise, public, versioned list draws sustained work from professionals, amateurs and,
increasingly, AI systems. Information theory has many famous open problems: the capacity of the
deletion channel, the two-receiver broadcast channel, the interference channel, the Shannon
capacity of $C_7$, the entropy region for four or more variables, and many more. These are
scattered across surveys, monographs and talk slides, with statements that differ in their
quantifiers and conventions.

We built a catalog of these problems as a git repository. Each problem has a Markdown file
(statement, background, what is known, references) and a Lean~4 file defining a proposition
\code{Statement}. The Lean files and the shared \code{Core} definitions are frozen: a SHA-256 lock
file covers 130 files, and CI rejects any solution that modifies them. A solution is a Lean file
proving \code{Statement} or its negation. A verifier checks definitional equality of the proved
type with the frozen proposition and that the proof depends only on \code{propext},
\code{Classical.choice} and \code{Quot.sound}, and replays the proof with \code{leanchecker}.

Freezing statements makes a solution's meaning unambiguous, but it also freezes any mistake. The
main lesson of this note is that such mistakes are common, even in statements written carefully
by people and agents who know the subject. They are also predictable in type, and most of them
can be caught mechanically before anyone tries to solve the problem.

\paragraph{Contributions.}
\begin{enumerate}[leftmargin=*,itemsep=2pt]
\item A catalog of 114 open problems across ten areas of information theory, each with a frozen
Lean~4 statement and a kernel-checked solution protocol (Section~\ref{sec:catalog}).
\item \emph{Certificates}: Lean proofs about the exact objects in every frozen statement, showing
that the statement is not vacuous and agrees with known special cases. Coverage is 114/114, built
on a library of formalized Shannon theory (Section~\ref{sec:cert}).
\item An empirical account of what went wrong. 24 of the 109 originally frozen statements were
semantically repaired, 2 were already theorems, and 1 catalog formulation was formally disproved.
We give a six-way taxonomy and the detection method for each case (Section~\ref{sec:errors}).
\item The first formal results against the frozen statements: a kernel-checked disproof, and a
bridge from an external Lean proof of the Courtade--Kumar conjecture (Section~\ref{sec:results}).
\item Recommendations for formal benchmarks of open problems (Section~\ref{sec:lessons}).
\end{enumerate}

\section{The catalog}\label{sec:catalog}

\begin{table}[t]
\centering
\begin{tabular}{lr@{\qquad}lr}
\toprule
Area & Problems & Area & Problems\\
\midrule
Network information theory & 18 & Channel capacity & 12\\
Coding theory & 13 & Quantum information & 12\\
Entropy inequalities & 13 & Network coding & 10\\
Source coding & 10 & Zero-error information theory & 10\\
Security and privacy & 8 & Statistical information theory & 8\\
\bottomrule
\end{tabular}
\caption{The 114 cataloged problems by area.}\label{tab:areas}
\end{table}

Table~\ref{tab:areas} lists the areas. Problems range from the classical (the binary deletion
channel, the general relay channel, the Shannon capacity of $C_7$, the four-variable entropy
region) to recent questions with precise conjectured answers (e.g.\ the Tahmasbi--Bloch
second-order conjecture for covert communication, or the minimum energy per bit of the Gaussian
relay channel under linear relaying).

\paragraph{Core definitions.}
All statements are written over a shared, frozen \code{Core} of about 2{,}000 lines. It contains
finite distributions, entropy and mutual information in bits, discrete memoryless channels and
their powers, operational capacity regions defined through codes with vanishing error, codes and
distances, graphs and zero-error capacities, network codes, quantum states and channels,
differential entropy, and a few security and statistics primitives. Operational quantities are
defined from codes, not from single-letter formulas. For example, the broadcast-channel capacity
region is the closure of the set of rate pairs achievable by sequences of codes, so a
single-letter formula cannot be built into a definition by accident.

\paragraph{Solutions and attempts.}
Formal solutions live in \code{Solutions/}. Informal work (partial results, numerics, surveys,
dead ends) is recorded in per-problem attempt files that follow a protocol. In that protocol
every claim is tagged \code{[proved]}, \code{[sketch]}, \code{[heuristic]} or
\code{[conjectural]}, and the file's status moves from \emph{unverified} to
\emph{community-reviewed}, \emph{verified} or \emph{refuted}. The repository currently holds 64
attempt files.

\section{Certificates}\label{sec:cert}

A frozen statement can be wrong in ways that type-checking never reveals. The commonest case in
Lean is a \emph{junk value}. \code{sSup} of an unbounded or empty set of reals is $0$, division
by zero is $0$, and a capacity defined with such an \code{sSup} may silently equal $0$. The
``open problem'' then becomes trivial or meaningless.

For each problem we therefore require a \emph{certificate}, a Lean module that proves facts about
the \emph{exact} objects named in \code{Statement}:
\begin{itemize}[leftmargin=*,itemsep=2pt]
\item \textbf{Known bounds.} Every fact the docstring claims (``it is known that $A\le X\le B$'')
is proved about the frozen $X$. The proof uses the defining sets, not a re-definition.
\item \textbf{Non-vacuity.} The hypotheses of the statement are satisfiable, and the sets inside
each \code{sSup}/\code{sInf} are nonempty and bounded where they need to be.
\item \textbf{Regressions.} The frozen objects reproduce textbook special cases. Examples: the
Han--Kobayashi region equals the strong-interference capacity region; Marton's region contains
the degraded-BC capacity region; decode-and-forward is at most the frozen relay inner bound; the
successive-cancellation error of an uncoded bit on BSC$(p)$ is $p$; the complementary graph
entropy of $K_n$ under the uniform law is $\log_2 n$.
\item \textbf{Settled statements.} If a certificate proves \code{Statement} or
\code{\textasciitilde Statement} outright, the problem was not open. The audit lists it
separately and excludes it from coverage, and the statement goes to a repair pass.
\end{itemize}
CI builds all certificates, audits every declaration for axioms beyond the standard three and for
name collisions, and replays each module with \code{leanchecker}. Coverage is currently 114/114,
with one statement listed as settled-false (Section~\ref{sec:results}).

\paragraph{The supporting library.}
The certificates needed real information theory, not just unfolding. Over the same period,
two cooperating agents built about 60{,}000 lines of ``bridge'' modules connecting \code{Core}
to Mathlib, plus about 10{,}000 lines of textbook results. Together they contain roughly 3{,}900
theorem and lemma declarations. Among them are chain rules, Fano, the AEP and typicality, Pinsker,
the method of types, Sanov, and Chernoff--Stein. They also include operational coding theorems
stated against \code{Core}:
\begin{itemize}[leftmargin=*,itemsep=1pt]
\item channel coding and the AWGN capacity $\tfrac12\log_2(1+P/N)$;
\item Slepian--Wolf, Wyner--Ziv and Gelfand--Pinsker;
\item the Gaussian MAC region;
\item Marton achievability (\code{martonInnerBound\_subset\_bcCapacityRegion});
\item the physically degraded relay capacity;
\item wiretap secrecy capacity;
\item the one-dimensional Shannon entropy power inequality via Stam and de~Bruijn;
\item minimax universal redundancy asymptotics;
\item Lov\'asz's $\Theta(C_5)=\sqrt5$.
\end{itemize}
Certificates are about 34{,}000 lines.

\section{What went wrong}\label{sec:errors}

The first frozen release had 109 problems (be80f07, 26~Nov~2026). Five problems added later are
excluded from the statistics below. We count a statement as \emph{repaired} if a later commit
changed the meaning of its Lean file. Docstring-only edits do not count. Two \code{Core}
definitions (the hypercontractivity ribbon and multi-source networks) were also corrected; the
statements that use them are not added to the count.

\begin{itemize}[leftmargin=*,itemsep=2pt]
\item \textbf{24 of 109 (22\%)} statements were semantically repaired. Six were repaired more
than once: Gaussian relay three times, universal lossy redundancy three times, and the Gaussian
interference channel, three-description Gaussian multiple descriptions, explicit list decoding and
the timing channel twice each.
\item \textbf{2} statements were found to be known theorems. The sharp R\'enyi entropy power
inequality statement is Ram--Sason (2016), Thm.~1, which holds for all independent pairs. The formalized matrix-SNR
single-crossing property of the MMSE follows from Payar\'o--Palomar (2009) and
Bustin--Payar\'o--Palomar--Shamai (2013). Both are now marked solved.
\item \textbf{1} statement, a formulation of the catalog's own that is not posed in the
literature, was formally disproved: the SDPI equality question (Section~\ref{sec:results}).
\item A provenance audit also flagged one further catalog formulation (LDP minimax risk under KL)
and one misattributed conjecture (batch codes). Neither changed the Lean statement.
\end{itemize}
In total, 27 of the 109 original statements (25\%) did not, as frozen, encode a genuinely open
problem.

\begin{table}[p]
\centering\footnotesize\sloppy
\begin{tabularx}{\textwidth}{>{\raggedright\arraybackslash}p{2.6cm}cL>{\raggedright\arraybackslash}p{2.3cm}}
\toprule
Type & \# & Examples (problem: defect) & Caught by\\
\midrule
T1. Lean artefacts (junk values, computability, encodings) & 5 &
AWGN feedback, Gaussian MD3, Gaussian relay: \code{sSup}/\code{sInf} in $\mathbb R$ over
possibly unbounded sets (moved to $\mathbb R_{\ge0}\cup\{\infty\}$ with guards); explicit list
decoding: a Turing-machine decoder had to compute $\lfloor(1-h(p)-\varepsilon)n\rfloor$ for
non-computable reals; earlier, its tape had no blank delimiter. LDP minimax: zero-mass
estimators gave infinite KL. &
Re-audit; non-vacuity certificates\\
T2. Mis-transcribed known regions or definitions & 7 (+2 \code{Core}) &
Relay: Cover--El Gamal Thm.~7 transcribed without its auxiliary structure. Gray--Wyner with side
information: bin index charged only to the common link. Three-receiver BC: a hand-made Marton
region. Gaussian IC: time sharing without power control or coded $Q$. BMC: the plain
Dueck/Berger bound, vacuous here. Polar codes: the ``SC decoder'' was any function, i.e.\
block-MAP decoding. Distributed hypothesis testing: the Shimokawa--Han--Amari exponent without time sharing. Hypercontractivity ribbon (\code{Core}): wrong orientation and form, replaced
by Nair's KL characterization. Network coding (\code{Core}): sources could equal sinks. &
Regression certificates; literature re-reading\\
T3. Statements false as posed (contradicted by known results or numerics) & 7 &
Gaussian relay: hybrid DF/CF optimality, false by El Gamal--Mohseni--Zahedi (2006); then a
low-SNR slope question contradicted by El Gamal--Gohari--Nair (2022). Primitive BSC relay:
$\max(\text{PDF},\text{CF})$, beaten by time sharing. Gaussian MD3: equality with a single
scheme. Poisson: $O(\sqrt A)$ support, against $A^{0.71}$ numerics. Timing: a relaxed constraint
made the statement false for every service law. Universal lossy: full-simplex minimax. Sperner:
$\Sigma(\vec C_5)<1$, but $\Sigma(\vec C_m)=1$ for every directed cycle. &
Literature; numerics; explicit counterexamples\\
T4. Already settled questions & 3 (+2) &
Covert communication: the keyed-vs-keyless question is settled by Tahmasbi--Bloch Thm.~1.
Wyner--Ziv dispersion: a single-dispersion form, replaced by a two-constant $(V_{\min},V_{\max})$ form (the reason is recorded only in the diff; counted here provisionally). Noisy
interactive coding: a non-adaptive model, against the adaptive randomized one in which the
question is posed. Plus the two known theorems above (R\'enyi EPI, MMSE). &
Literature search\\
T5. Unposed or misattributed questions & 2 (+1) &
Zero-error BC: computability of a feedback capacity (posed nowhere), retargeted to Hu--Shayevitz's
$\rho^*(C_5)>0$. Graph entropy: the open question concerns complementary graph entropy, not
Witsenhausen's rate. Batch codes: misattributed conjecture (no change to the Lean statement). &
Provenance audit (verbatim quote required)\\
T6. Degenerate quantifier domains & 2 &
Network-coding decidability: networks with only local demands. Cognitive IC: a missing rate
constraint in the region. &
Re-audit\\
\bottomrule
\end{tabularx}
\caption{Taxonomy of defects in the 109 originally frozen statements. Counts are by problem and
type. A problem repaired twice for different reasons is counted under both types, so the column
sums to more than 24. Details and diffs are in commits 0c19d64, 094f988, 49c683f, 36fb580,
03ced53, 440e082, 6e9452a, a6a7be1, 04eb1a9, b0d03d1, bcfe4c9, d2acd3a and f8059b4.}
\label{tab:taxonomy}
\end{table}

Table~\ref{tab:taxonomy} gives the taxonomy. Three observations stand out.

\paragraph{The hardest errors are faithful-looking transcriptions.} Type T2 errors pass every
syntactic check and most human readings. The relay-channel statement transcribed Cover--El
Gamal's Theorem~7 in a natural but wrong way, and an explicit $4\times2\to6\times2$ counterexample
refuted the transcription. The Gray--Wyner inner bound charged the bin index of the common
description only to the common link. A Slepian--Wolf regression on a BEC source exposed it:
the point $(e,0,1-e)$ is achievable at $D=0$ but was outside the frozen region. These were found
by \emph{regression certificates}: proving that the frozen objects reproduce a textbook special
case.

\paragraph{``Open'' needs a citation.} Four statements (the two T5 statements and two catalog
formulations, SDPI and LDP minimax) had no verifiable source posing them in the frozen form. Requiring a verbatim quote
from a source that poses the question was the single most effective provenance check.

\paragraph{Errors cluster in network information theory and in operational definitions.} Of the
24 repaired statements, 11 are in network information theory or multiterminal source coding.
These are the areas with the most auxiliary random variables, time-sharing conventions and
closure operations.

\section{Formal results}\label{sec:results}

\paragraph{A kernel-checked disproof (SDPI).}
The statement \code{SdpiKlVsChisqEquality} asked whether equality of the KL and $\chi^2$
strong data-processing constants on both sides forces membership in a hypercontractivity ribbon.
It is a formulation of the catalog's own. A three-state reversible counterexample was first found
numerically; the solution file (\code{Solutions/SdpiKlVsChisqEquality.lean}) then proves
$\neg$\code{Statement} with standard axioms only. The key step is a global inequality
$D(PW\|QW)\le 0.09\,D(P\|Q)$ that is tight to second order at $P=Q$. It is closed by a grouping
identity, a power-series bound, and an exact-rational LP positivity certificate. The value of
this result is mainly as a catalog correction. It also shows how cheaply a mis-specified
statement falls to a determined search.

\paragraph{A second disproof (finite-letter exact common information).}
The solution file for \code{CommonInformation\-ExactVsWyner} disproves finite-letter attainment
of exact common information. On the symmetric binary erasure source SBES$(3/4)$, every finite-letter
value strictly exceeds $h(3/4)$ while the limit is at most $h(3/4)$. The proof is about 2{,}600
lines and combines Pinsker-based strictness with a shell-wise random-rectangle covering bound. Its
mathematical content merits a separate write-up.

\paragraph{A bridge to an external proof (Courtade--Kumar).}
The Courtade--Kumar conjecture ($I(f(X);Y)\le1-h(\alpha)$ for Boolean $f$ under BSC$(\alpha)$
noise) was proved in September 2026 by several groups. One proof (Chen et~al., arXiv:2609.24931)
is formalized in Lean in a separate repository of about 45{,}500 modules on Lean v4.33.1. The
catalog uses v4.36.0, so the proof cannot be imported. Its target proposition, however, depends
only on Mathlib.

We copied the target file byte-for-byte (pinned commit and SHA-256) and proved in the catalog
\[
\code{GeneralCK.GeneralCourtadeKumar}\;\Longrightarrow\;\code{CourtadeKumar.Statement}.
\]
The core lemma is an \emph{equality}: the catalog's $I(f(X);Y)$, a pushforward of a uniform
input through a product BSC, and the external project's hand-written finite-sum formula agree for
every $n$, every $f$ and every $\alpha\in[0,1]$. Two formalizations written independently,
by different authors and for different purposes, agree exactly; this is itself evidence that neither is
mis-specified. The chain to the frozen statement has two kernel-checked hops in separate builds.
A single-kernel check would need a bridge project on the external toolchain.

\paragraph{Claimed resolutions are tracked, not trusted.}
Several catalog problems are claimed resolved by preprints from 2026. Examples are Marton's inner
bound (Huang--Liu--Liu, arXiv:2608.19869), the Li--M\'edard conjecture (Yu, arXiv:2611.00868),
sharp continuity of quantum conditional entropy (arXiv:2607.24687), and good qLDPC codes. The
catalog records these as \code{claimed-resolved}, with the evidence and the gap. For example,
Huang--Liu--Liu's explicit numbers have been independently recomputed, but their certified
one-letter upper bound has not. The status changes only once a check is recorded.

\FloatBarrier
\section{Recommendations}\label{sec:lessons}

\begin{enumerate}[leftmargin=*,itemsep=3pt]
\item \textbf{Certify before you freeze.} Prove the known bounds about the exact frozen objects,
and prove non-vacuity of every hypothesis and every \code{sSup}/\code{sInf}. Junk values are the
first thing a prover will exploit.
\item \textbf{Write regression certificates against textbook special cases.} Most type T2 errors
were caught this way, and nothing else caught them reliably.
\item \textbf{Prefer operational definitions.} Defining capacities through codes, not formulas,
kept every single-letter formula a theorem to be proved rather than a definition.
\item \textbf{Require a verbatim source for ``open''.} Eight of the 27 defective statements
were either posed nowhere in the frozen form or already settled.
\item \textbf{Search for settled-ness, and re-search.} The literature moves quickly. Two
statements were already theorems, and four more problems saw claimed resolutions within months.
\item \textbf{Make retirement explicit.} When a statement is repaired, keep the old proposition,
renamed and marked retired, together with its disproof. Never keep it as the problem.
\item \textbf{Cross-check against independent formalizations.} When an outside group formalizes
the same quantity, proving the two definitions equal is cheap. It is also strong evidence for both.
\end{enumerate}

\section{Limitations}

Certificates establish non-vacuity and agreement with special cases. They cannot establish that a
statement says what the literature means. Most of the auditing described here was done by AI
agents working under an explicit protocol, with limited expert human review, so some of the 82
statements not flagged may still be wrong. The defect counts cover a short, intensive period and
should be read as a lower bound on the error rate of carefully written, unaudited formal
statements. The counts by type involve judgment calls; the commit history gives the raw material.

\section*{Acknowledgements}
We thank the authors of the external Lean development \code{dpwoodru/general-courtade-kumar-lean}
for making their target statement self-contained, which made the bridge in
Section~\ref{sec:results} a short proof.

\begin{thebibliography}{99}\small
\bibitem{cgjlmnw} Z.~Chen, A.~Gohari, A.~Javanmard, H.~Lin, V.~Mirrokni, C.~Nair, D.~P.~Woodruff,
``A proof of the most informative Boolean function conjecture,'' arXiv:2609.24931, 2026.
\bibitem{hll} M.~Huang, Y.~Liu, Y.~Liu, ``Sub-optimality of Marton's inner bound for the
two-receiver broadcast channel,'' arXiv:2608.19869, 2026.
\bibitem{yu} L.~Yu, ``Noise stability on the hypercube: order $7/8$,'' arXiv:2611.00868, 2026.
\bibitem{rs} E.~Ram, I.~Sason, ``On R\'enyi entropy power inequalities,'' \emph{IEEE Trans.\
Inf.\ Theory}, 2016 (arXiv:1601.06555).
\bibitem{emz} A.~El Gamal, M.~Mohseni, S.~Zahedi, ``Bounds on capacity and minimum
energy-per-bit for AWGN relay channels,'' \emph{IEEE Trans.\ Inf.\ Theory}, 2006.
\bibitem{egn} A.~El Gamal, A.~Gohari, C.~Nair, ``Achievable rates for the relay channel with
orthogonal receiver components,'' IEEE ITW, 2021; and ``A strengthened cutset upper bound on
the capacity of the relay channel and applications,'' \emph{IEEE Trans.\ Inf.\ Theory}, 2022.
\bibitem{tb} M.~Tahmasbi, M.~R.~Bloch, ``First- and second-order asymptotics in covert
communication,'' \emph{IEEE Trans.\ Inf.\ Theory}, 2019.
\bibitem{hs} Y.~Hu, O.~Shayevitz, ``The $\rho$-capacity of a graph,'' \emph{IEEE Trans.\ Inf.\
Theory}, 2017.
\bibitem{nair} C.~Nair, ``Equivalent formulations of hypercontractivity using information
measures,'' IZS, 2014.
\bibitem{mathlib} The mathlib Community, ``The Lean mathematical library,'' CPP, 2020.
\end{thebibliography}
\end{document}
